\documentclass[12pt]{article}
\usepackage[utf8]{inputenc}
\usepackage[T1]{fontenc}
\usepackage{hyperref}
\title{Tambotalo Language Resources for K-12 Curriculum Development}
\author{Compiled by Claude Sonnet 5 (claude-sonnet-5) for Andr\'as Kornai}
\date{}
\begin{document}
\maketitle

\section{Introduction}
Tambotalo (ISO 639-3 tls; Glottocode tamb1253), also called Biliru, is an Oceanic language spoken in one village in southeast Espiritu Santo, Vanuatu. Search results rate it severely endangered. A 1981 count of about 50 speakers is cited, with fewer likely today. Neighbouring languages (Tutuba, Naranggo, Shark Bay) have drawn speakers away. It has no literacy tradition.

\section{Primary Educational Resources}
None found.

\section{Linguistic and Reference Materials}
\begin{itemize}
\item Wikipedia, Tambotalo language: \url{https://en.wikipedia.org/wiki/Tambotalo_language}
\item Austronesian Basic Vocabulary Database entry: \url{https://abvd.eva.mpg.de/austronesian/language.php?id=1084}
\item Glottolog languoid page: \url{https://glottolog.org/resource/languoid/id/tamb1253.bigmap.html}
\end{itemize}

\section{Religious and Community Resources}
Joshua Project lists the people group with a very small or unknown population and says it may have merged into surrounding communities: \url{https://joshuaproject.net/people_groups/15231}

\section{Digital and Audio Resources}
\begin{itemize}
\item Endangered Languages Project entry: \url{https://endangeredlanguages.com/lang/2509}
\end{itemize}

\section{Curriculum Development Context}
Beyond a wordlist and catalogue entries, I found no documentation. A school curriculum is not realistic now. Documentation work with the remaining elderly speakers would have to come first.

\end{document}